\chapter{Technický abstrakt}
\label{chapter:Technical-abstract}
\pagenumbering{arabic}

Empirické zistenia o vplyve nástrojov umelej inteligencie na učenie si zdanlivo protirečia: adaptívny tutor prekonáva aktívne učenie~\cite{kestin2025aitutoring}, generatívny nástroj však znižuje výsledky na následnom teste~\cite{bastani2025harm}, a porovnanie adaptívneho a generického nástroja nenachádza rozdiel~\cite{slijepcevic2025khanmigo}. Práca tvrdí, že rozpor je zdanlivý a zjednocuje ho jedna hypotéza: nástroj zvyšuje výkon počas hodiny a pocit naučenia takmer vždy, no jeho vplyv na \emph{trvalé} osvojenie závisí od toho, či interakcia žiaka s ním zostáva generatívna, alebo sa stáva offloadingová. Táto veličina však dosiaľ nebola priamo meraná~--- Gerlich~\cite{gerlich2025offloading} ju získava sebahodnotením, Kosmyna~\cite{kosmyna2025brain} elektroencefalografiou nepoužiteľnou v triede.

Práca predkladá dva prepojené príspevky. \textbf{Prvým} je \emph{index epistemickej kvality promptov} (PEQI): taxonómia žiackych dopytov ukotvená v rámci ICAP~\cite{chi2014icap}, ktorá radí prompty od žiadosti o hotové riešenie po predloženie vlastnej úvahy na posúdenie, spolu s automatickým klasifikátorom postaveným na priehľadných lexikálnych príznakoch. Na kotviacom súbore 90~slovenských promptov dosahuje klasifikátor v krížovej validácii $\kappa=0{,}92$, pričom 98,8\,\% chýb je posunom o jedinú susednú úroveň. Index sa počíta z promptov zachytených nástrojom vyvinutým v technickej časti práce; jeho prvá iterácia postavená na knižnici LittleProxy zlyhala na nepodpore HTTP/2 a WebSocket, a bola nahradená architektúrou s vymeniteľnou záchytnou vrstvou nad mitmproxy.

\textbf{Druhým} príspevkom je plne špecifikovaný krížový experiment: replikovaný latinský štvorec $3\times3$ so šiestimi sekvenciami, v ktorom každý žiak absolvuje klasickú výučbu, výučbu s generickým nástrojom aj s adaptívnym sokratovským nástrojom na troch izomorfných témach algoritmizácie. Generická a adaptívna podmienka používajú \emph{ten istý model, tie isté úlohy a to isté rozhranie}; manipulovaná je výhradne pedagogická vrstva, čím sa izoluje účinok ohraničenia od účinku modelu~--- čo je zámena, ktorá obmedzuje doterajšiu literatúru. Okamžitá a odložená forma testu sú prepojené kotviacimi položkami a kalibrované súbežne, takže veľkosť zabúdania je merateľná a nie iba predpokladaná.

Osobitná pozornosť je venovaná tomu, čo randomizácia unesie. Prideľovaná je podmienka, nie kvalita promptov, takže mediačné tvrdenie stojí na predpoklade, ktorý dizajn nedodáva. Do konfirmačnej rodiny preto vstupujú len tri tvrdenia, ktoré sú identifikované randomizáciou a môžu zlyhať; mechanizmus sa testuje trianguláciou troch odhadov opierajúcich sa o tri rôzne predpoklady~--- mediačný model, randomizovaná pedagogická vrstva ako inštrumentálna premenná, a dávková závislosť v rámci jednej podmienky~--- doplnenou o analýzu citlivosti, ktorá vyčísli, aké silné by nepozorované zmätenie muselo byť na zvrátenie zistenia.

Druhá vrstva tej istej otázky sa týka merania. Mediátor nie je pozorovaný, ale odhadnutý klasifikátorom, a chyba toho odhadu môže závisieť od randomizovanej podmienky: adaptívna vrstva žiakom predpisuje jazykovú formu, ktorú klasifikátor číta ako doklad uvažovania. Práca ukazuje, že prvá verzia klasifikátora bola na túto zámenu plne citlivá~--- štyri slová bez obsahu posunuli každý z \venN{} kotviacich promptov najnižších úrovní na najvyššiu, čo je $\venDeltaBefore$ indexu na prompt~--- a odvodzuje, čo by taká chyba spôsobila: redukovaná forma je voči nej odolná, inštrumentálny odhad utlmený faktorom $a/(a+\delta)$ a mediačný odhad skreslený v oboch smeroch tak, že sa skreslenia môžu presne zrušiť. Z toho plynie poradie vykazovania, kalibrované očakávanie pre trianguláciu a zaslepený stratifikovaný audit dimenzovaný na silu $\audPower$.

Analytický plán je uzavretý, implementovaný a spustený na dátach zo známeho generujúceho modelu pred zberom. Validácia potvrdila obnovenie vložených účinkov a odhalila šesť chýb návrhu, ktoré by sa inak prejavili až po zbere: singularitu v kódovaní prenosu medzi periódami, nedostatočnú dĺžku odloženého testu, nedostatočnú veľkosť vzorky pre tvrdenie o ekvivalencii, prítomnosť neidentifikovateľnej hypotézy v konfirmačnej rodine, chýbajúce prepojenie testových foriem a chýbajúce kritérium presnosti pre mechanizmus. Cieľová veľkosť vzorky sa tým zvýšila z 12 na 18 tried, pričom k tejto hodnote vedú dve nezávislé obmedzenia: test ekvivalencie a požiadavka na šírku intervalu pre mechanizmus.
